\documentclass[11pt]{article}
\usepackage[margin=1in]{geometry}
\usepackage{amsmath,amssymb,amsthm,booktabs}
\usepackage[T1]{fontenc}
\usepackage{hyperref}
\newtheorem{theorem}{Theorem}
\newtheorem{lemma}[theorem]{Lemma}
\newtheorem{corollary}[theorem]{Corollary}
\newcommand{\Con}{\operatorname{Con}}
\newcommand{\Halt}{\operatorname{Halt}}
\newcommand{\Native}{\mathsf{NativeAx}}
\newcommand{\MP}{\mathsf{MP}}
\title{Semantic targets and a tag-independent inference rule\\
for the existing consistency-search calculus}
\author{Research notes prepared for Vladimir Reshetnikov's ProveIt project}
\date{18 October 2026}
\begin{document}
\maketitle

\begin{abstract}
Two changes to the existing proof-search source are justified without changing
its set theory or its formula codec. First, the halting test may recognize every
membership atom, rather than only the distinguished self-membership atom.
Second, modus ponens may ignore the outer constructor tag of its cached premise,
provided the two encoded fields still match the current theorem and the proposed
conclusion. The latter change is sound specifically at theorem boundaries;
ordinary quantifier elimination handles the universal tag, and the native
container and irreflexivity axioms handle the equality tag. The existing concrete
proof-producing command blocks remain complete. These are ordinary mathematical
source-level results, with finite semantic controls, not measured state savings,
new binary tables, checked Lean theorems, or a completed PA internalization.
\end{abstract}

\section{Exact setting and relationship to prior work}

The inspected repository is ProveIt at commit
\texttt{e58338db7545909907e07a7603e8eb88625c527a}.
The combined243 baseline uses the original complete numerical calculus, with
native equality, the CP/GIC bundle, and both freshness gates. Its ordinary
ZFC consistency calibration is retained separately from its relative Lean
machine comparison and the unfinished internal arithmetic endpoint
\cite{combined,native}.

Let
\[
 P(a,b)=\frac{(a+b)(a+b+1)}2+b.
\]
The formula constructors are
\[
 \begin{array}{ll}
 I(a,b)=4P(a,b),& R(i,j)=4P(i,j)+1,\\
 Q(i,a)=4P(i,a)+2,& E(i,j)=4P(i,j)+3.
 \end{array}
\]
They mean implication, membership, universal quantification, and equality,
respectively. The distinguished contradiction is
\(F=1=R(0,0)\). Code zero has the stipulated truth meaning; it is not parsed
by infinitely unfolding the equation \(0=I(0,0)\).
All positive codes have their actual finite decoded meaning, including code
fields that serve as variable names.

The native first-order theory consists of the transitive power-bounding
container sentence CP, the full GIC schema with arbitrary finite parameters
and logical equality, and irreflexivity
\(\forall x\,\neg(x\in x)\). Its axioms are sentences. In particular CP entails
the elementary hosting consequence
\[
 \forall x\,\exists y\;(x\in y).                    \tag{Host}
\]
The original compiler and proof invariant are attributed to Mira-Cult, Jeff
Ketchersid and contributors, Andrew James Wade, and the maintained ProveIt
development. The present notes do not replace their source calibration.

Earlier repository results include changes to successful-MP result placement,
dummy-stack elimination, shared comparisons, a specialized test for the
single cache value one, and decoding a cached implication before testing its
antecedent \cite{cachemp,dummy,decoded}. The two changes below instead use
semantic restrictions on which \emph{whole formulas can be theorems}. They were
not found in the inspected relevant notes. This is a scoped comparison with the
repository, not a claim of historical priority in proof theory or compilation.

\section{Every atomic membership theorem is a refutation}

Write \(M\models_* A\) when a decoded formula \(A\) is true in \(M\) under every
assignment of its free names. The distinction between this condition and truth
at one chosen assignment is essential.

\begin{lemma}[Atomic membership]
Let \(M\) be a nonempty irreflexive membership structure. For all variable
indices \(i,j\), \(M\not\models_* R(i,j)\).
\end{lemma}
\begin{proof}
Choose any object \(u\) of the domain and give every variable name the value
\(u\). The atom becomes \(u\in u\), which is false. This includes both
\(i=j\) and \(i\ne j\); no empty-set axiom, extensionality, container axiom,
Foundation, or Choice is needed.
\end{proof}

Consequently a derivation of \(R(i,j)\) from a sentence theory containing
irreflexivity gives an ordinary finite refutation. A direct syntactic route
generalizes the free names, instantiates them by one common fresh name, and
combines the resulting diagonal membership atom with irreflexivity. Coincident
indices are handled by one generalization instead of two. Renaming every symbol
of an arbitrary proof to one name is not a sound substitute for this argument:
it may capture variables inside the proof.

Define the numerical target set
\[
 \mathcal B=\{h\in\mathbb N:h\equiv1\pmod4\}.
\]
It contains \(F\), and every member decodes to one of the atoms covered by the
lemma. Keep every command unchanged except that NIP halts when its incoming
cache lies in \(\mathcal B\). On all other inputs it performs the original
cache clear and stack pop.

\begin{theorem}[Expanded NIP target]
From an initialization satisfying the original theorem invariant and its
original completeness interface, the changed source enumerator halts exactly
when the original source enumerator halts. In particular it has the same
ordinary ZFC consistency calibration.
\end{theorem}
\begin{proof}
Before the first newly accepted halt, the two computations coincide literally
at command boundaries. If the old computation halts, the new computation
either has already halted or reaches the same cache value \(F\) and halts.
Conversely, at any new halt the cached formula is a theorem by the inherited
invariant. The lemma turns that atomic theorem into a refutation; original
proof-search completeness therefore gives halting of the old enumerator.
All arguments concern finite traces and finite proofs and can be carried out
inside ZFC. No consistency assumption is needed as a hypothesis of the
equivalence theorem.
\end{proof}

This is not equality of NIP effects on arbitrary register files: for example,
cache value five causes an immediate new halt and an ordinary old NIP step.
Likewise the statement does not justify changing arbitrary binary rows merely
because their intended register names look related. A concrete changed machine
would need its own runtime proof and a counted blank initialization.

\subsection{A complete scalar implementation, with honest costs}

The cache can be consumed by a four-test cycle:
\begin{quote}\ttfamily
repeat:\\
\quad if H=0: go to POP; otherwise decrement H\\
\quad if H=0: HALT; otherwise decrement H\\
\quad if H=0: go to POP; otherwise decrement H\\
\quad if H=0: go to POP; otherwise decrement H\\
\quad return to repeat.
\end{quote}
Every nonhalting exit has \(H=0\), so the original pop can follow with its
original scratch assumptions. The loop performs exactly \(h\) successful
decrements and one failing zero test. Its control has four residue positions
and needs no arithmetic scratch register. The current specialized NIP call has
an increment of scratch1, a comparison wrapper, and a halt wrapper, while the
generic comparison and clear loops are shared with other commands. Therefore
those shared states cannot all be deducted from the machine budget. Four
residue controls may cost more than the present singleton check. No binary
state improvement is claimed from this implementation.

\section{Ignoring the cached outer tag in modus ponens}

The original MP command, with proposed conclusion \(b=d+x\), tests
\(H=I(T,b)\). Replace that test by
\[
       \left\lfloor\frac H4\right\rfloor=P(T,b).        \tag{TMP}
\]
Keep the original success and failure effects, including the dummy truth push
on success, the additive parameter, and the cache clear. The condition accepts
the four possible formulas \(I(T,b),R(T,b),Q(T,b),E(T,b)\).

\begin{lemma}[A nontrivial domain]
Every nonempty model of CP and irreflexivity contains two distinct objects.
\end{lemma}
\begin{proof}
Choose \(u\). Hosting gives \(v\) with \(u\in v\). If \(u=v\), irreflexivity
would be violated. Thus \(u\ne v\).
\end{proof}

\begin{theorem}[Tag-independent theorem inference]
Let \(M\) be an irreflexive membership structure with at least two objects.
If \(M\models_* T\), \(M\models_* H\), and (TMP) holds, then
\(M\models_* b\).
\end{theorem}
\begin{proof}
Use Euclidean division to write \(H=4P(T,b)+r\), where \(0\le r<4\).
\begin{itemize}
\item If \(r=0\) and \(H>0\), the cached formula is \(T\to b\), so ordinary
modus ponens gives the assertion. If \(H=0\), then \(P(T,b)=0\) forces
\(T=b=0\), and the conclusion is the stipulated truth code.
\item If \(r=1\), the cache is the atomic membership formula
\(v_T\in v_b\). The preceding atomic lemma contradicts its validity under
every assignment. Thus this case cannot occur for the stated semantic inputs.
\item If \(r=2\), the cache is \(\forall v_T\,b\). For an arbitrary
assignment \(e\), instantiate its universal quantifier at \(e(T)\); the
updated assignment is \(e\) itself. Hence \(b\) is true at every assignment.
There is no freshness restriction between the index \(T\) and the formula
code \(b\).
\item If \(r=3\), the cache is \(v_T=v_b\). If the numerical indices differ,
assign two distinct domain objects to these two names. This contradicts
validity under all assignments. Therefore \(T=b\) as natural numbers, and
the desired formula is the already valid current formula.
\end{itemize}
\end{proof}

This is a theorem-level admissible inference, not the assertion that the
formula \((T\land E(T,b))\to b\) is a tautology. At one assignment distinct
variable names may have equal values. For example take current truth code
\(T=0\), proposed contradiction \(b=1\), and cache \(E(0,1)=11\). Under an
assignment giving names zero and one the same value, both premises are true
and the conclusion is false in an irreflexive structure. Only the stronger
all-assignment premise rules this out.

The two-object premise is also indispensable. In the singleton structure with
empty membership, both code zero and \(E(0,1)=11\) are valid, while \(F=1\)
is not. That structure is excluded by CP. Thus the equality case cannot be
justified by citing irreflexivity alone or by dropping the source's set theory.

\subsection{The existing command compiler remains complete}

The original source has a concrete compiler from numerical derivations to
command blocks. The important issue is that a compiler can intentionally use
a \emph{failed} MP to clear a register; merely accepting every old successful
MP would not prove preservation of that compiler.

The exact inspected \texttt{Enum254.lean} defines
\[
 \mathrm{discard}=[\MP(1),\mathrm{AX}(1,0,0,0)]
                         \mathbin{+\!+}\mathrm{NIP}(0)^{10}.
\]
Its interface starts with cache zero and additive parameter zero. The first
MP still fails under (TMP), for
\[
 \lfloor0/4\rfloor=0<P(T,1).
\]
It therefore has exactly the old clearing and stack effect, for every initial
current formula \(T\) and every stack. AX and the following pops are unchanged.
Every MP in the remaining inductive proof blocks is an intended successful
comparison against the exact implication \(I(A,B)\), so it is still accepted
with exactly the original result and stack effect. The leaf, generalization,
and MP blocks consequently retain their original strong contracts, by the
same structural induction. Initial AX still clears the arbitrary additive
parameter and cache before a proof block. Command packing, pointer descent,
candidate enumeration, and next-candidate preservation have not changed.

\begin{corollary}[A complete alternative source]
Replace only the MP comparison by (TMP), leaving NIP's original test unchanged.
From the calibrated source initialization, the new enumerator halts if and
only if \(\Native\) is inconsistent, and hence has the inherited ordinary
ZFC consistency calibration.
\end{corollary}
\begin{proof}
If \(\Native\) is consistent, ordinary first-order completeness supplies a
nonempty set-coded model. CP makes it nontrivial. Maintain validity under all
assignments for the current formula, cached formula, and every stack element.
The theorem proves the only changed invariant step. A NIP halt would make
\(F\) valid, contradicting irreflexivity.

If \(\Native\) is inconsistent, the old calibrated numerical calculus has a
finite refutation. The preceding inspection of its concrete proof compiler
shows that the same blocks still construct that refutation in the changed
source, unless an earlier legitimate halt occurs. The unchanged fair candidate
enumeration reaches a halting block. These two directions are ordinary proofs
inside ZFC; they do not require external well-foundedness of a model.
\end{proof}

The rule can also be combined with the expanded NIP target. The extra early
halts remain justified, and the old compiler's discarded AX rows have outer
implication or universal tags, never membership tags.

For concrete implementation one can compare \(P(T,b)\) with a quotient of the
consumed cache, or compare \(4P(T,b)\) with the cache and accept residuals zero
through three. The former removes one formula-scaling occurrence but adds
quotient control; the latter extends a shared comparison and may require a
separate caller. Neither representation is free, and no compiler result or
new binary transition table is supplied here.

\subsection{A more concrete implementation: remove decoded-MP tag rejection}

Source57 already proposes a fully calibrated alternative that decodes a cached
implication, compares its antecedent with the current theorem, and returns its
decoded consequent without a dummy push \cite{decoded}. It explicitly rejects
the three non-implication tags. Its quotient is computed by consuming the cache
in groups of four and running the existing Cantor inverse-walk body once per
complete group. In that design exhaustion at the first decrement position
means a valid tag; exhaustion at any other position takes an invalid-tag
cleanup path.

The new inference theorem permits \emph{all four exhaustion positions to take
the same normal exit}. Let \((a,b)=P^{-1}(\lfloor H/4\rfloor)\). Return the
decoded consequent \(b\) if \(a=T\), otherwise return truth zero, preserving
the source's additive parameter. The group-drain loop consumes any incomplete
last group without invoking another inverse-walk body, so its output is exactly
the inverse pair of the quotient, for every cache value. No residue needs to
be stored. The original source57 moves the current theorem into the emptied
cache, compares it with the decoded antecedent, and transfers the saved
consequent to the current slot on equality; those operations still apply.

Soundness is precisely the tag-independent theorem on the matching branch,
and truth on the other branch. Source57's existing proof-producing blocks
remain complete: their successful calls supply genuine implications, while
normalization with cache zero returns current zero in both its matching and
nonmatching cases and leaves the underlying stack unchanged. This credits
the decoded-MP representation and its proof-block repair to the earlier work.
The new contribution is the justification for deleting its invalid-tag branch.
Whether that branch deletion frees an actual binary control state requires
compilation, runtime checking, and a new counted initializer; none is supplied
by this note.

\section{An exact restriction on odd theorem codes}

The same argument gives a useful semantic invariant stronger than the single
forbidden residue class. In every irreflexive structure with at least two
objects, an odd formula code is valid under every assignment precisely when
it is a reflexive equality atom. In the exact numerical codec this means
\[
 n\text{ odd and }M\models_* n
 \quad\Longleftrightarrow\quad
 \exists i\in\mathbb N\;n=4P(i,i)+3=8i(i+1)+3.
\]
In particular every valid odd code is congruent to three modulo sixteen;
that residue condition alone is not sufficient. Applied to the source theorem
invariant under consistency, this describes \emph{all} possible odd theorem
codes. Without a consistency assumption, any derivation of an odd code outside
the displayed family gives a refutation instead. A future minimizer could use
this arithmetic invariant, but proving that a concrete control state contains
a theorem code remains part of its correctness obligation.

\section{Effective proofs, internalization, and finite evidence}

The mathematical proofs use only finite syntactic manipulations, ordinary
first-order soundness/completeness on set-coded structures, and finite source
traces. They establish ZFC-formalizable source results. In a conventional
effective Hilbert or sequent calculus, the displayed direct derivations for
membership, universal elimination, and off-diagonal equality provide
primitive-recursive transformations of supplied \emph{native first-order}
proofs. The equality case uses one fixed native proof that two distinct objects
exist, then generalizes and instantiates the two distinct names.

No implemented arithmetic translator from the exact C254 proof codec through
those native proofs and back is provided. The repository's semantic
completeness bridges alone are not an implemented primitive-recursive
translator. In particular these notes do not supply a PA proof of the complete
new machine equivalence, a serialized internal ZFC derivation, or a new kernel
proof. A concrete binary implementation must also bind its exact table,
positive-progress runtime simulation, and complete blank initialization.

The independent program \texttt{semantic\_targets.py} uses exact integers and
the stated zero alias. Its completed controls record:
\begin{center}
\begin{tabular}{lr}
\toprule
Control&Cases\\
\midrule
Cantor round trips&10,000\\
Atomic targets in nonempty irreflexive structures&4,416\\
Odd-code validity characterization&8,704\\
Tag-independent MP instances&156,672\\
Of these, instances with valid premises&17,204\\
Complete scalar residue-drain controls&10,000\\
\bottomrule
\end{tabular}
\end{center}
The model checks exhaust irreflexive relations on the stated carriers of sizes
one, two, or three, and quantify over every assignment of the finitely many
names occurring in each tested formula. They are controls for the semantic
lemmas, not finite models of CP/GIC and not evidence for consistency strength.
The singleton and single-assignment counterexamples ensure the two necessary
hypotheses are actually tested. Nontrivial accepted new MP examples are
\((T,b,H)=(0,3,38)\) and \((3,3,99)\); both produce the theorem code three,
where the old exact-implication comparison would return zero.

\section{Alternative theory routes: what the audit supports}

The repository already has ordinary calibrations of raw equality-free PSC,
two-parameter EPSC, EPSC1 with exact BinaryUnion, finite assignment/class
kernels, and the equational encoding of finite universal clauses
\cite{raw,unary,equational}. These results should not be presented as newly
discovered routes. The one-parameter schema alone remains open: known support
models block raw Pairing, even with finite common transitive bounds, while
their internally pure parts prevent interpreting that separation as a
consistency-strength lower bound.

Existing complete theory controls do not justify a smaller-state claim. For
example the retained functional subset-closure experiment costs 347 working
states in its measured backend, against its 320-state subset-bound control;
its foundational calibration is simpler than its implementation
\cite{functional}. The recent finite equational packet has 2,324 bits before
its contextual-rewriting interpreter and initializer, and its smaller displayed
axiom count does not imply a smaller machine \cite{equational}.
The present two source results avoid these costs by retaining the calibrated
theory, all logical rules, and the existing formula code. Their actual state
cost is nevertheless an open engineering question.

\begin{thebibliography}{9}
\bibitem{combined} ProveIt, \emph{Combined 243-state candidate}, pinned repository
\href{https://github.com/VladimirReshetnikov/ProveIt/blob/e58338db7545909907e07a7603e8eb88625c527a/Computability/BusyBeaver/Research/ZFC/candidate-combined243/README.md}{README}.
\bibitem{native} ProveIt, \emph{An explicit first-order theory for the structured
program}, \href{https://github.com/VladimirReshetnikov/ProveIt/blob/e58338db7545909907e07a7603e8eb88625c527a/Computability/BusyBeaver/Research/ZFC/native_theory254.md}{native theory interface}.
\bibitem{cachemp} ProveIt closure-kernel intake, \emph{Source54: cache-result
modus ponens}, \href{https://github.com/VladimirReshetnikov/ProveIt/blob/e58338db7545909907e07a7603e8eb88625c527a/SetTheory/Cardinals/docs/reports/set-theory/zfc-closure-kernels/intake-158-source-54.md}{source54 review}.
\bibitem{dummy} ProveIt, \emph{Removing the successful-MP dummy push},
\href{https://github.com/VladimirReshetnikov/ProveIt/blob/e58338db7545909907e07a7603e8eb88625c527a/Computability/BusyBeaver/Research/ZFC/mp_dummy_push_partial_cost.md}{measured control}.
\bibitem{raw} ProveIt, \emph{Raw PSC: hereditary-saturation interpretation and
formalization}, \href{https://github.com/VladimirReshetnikov/ProveIt/blob/e58338db7545909907e07a7603e8eb88625c527a/Computability/BusyBeaver/Research/ZFC/raw_psc.md}{calibration and scope}.
\bibitem{unary} ProveIt, \emph{EPSC1 and an internally well-founded interpretation:
the remaining gap}, \href{https://github.com/VladimirReshetnikov/ProveIt/blob/e58338db7545909907e07a7603e8eb88625c527a/Computability/BusyBeaver/Research/ZFC/epsc1_wf_frontier.md}{one-parameter frontier}.
\bibitem{equational} ProveIt, \emph{Batch166: a compensated 244-state reduction},
\href{https://github.com/VladimirReshetnikov/ProveIt/blob/e58338db7545909907e07a7603e8eb88625c527a/SetTheory/Cardinals/docs/reports/set-theory/zfc-closure-kernels/intake-166.md}{equational branch and measured packet}.
\bibitem{functional} ProveIt, \emph{A functional subset-closure alternative: a
negative compiler result}, \href{https://github.com/VladimirReshetnikov/ProveIt/blob/e58338db7545909907e07a7603e8eb88625c527a/Computability/BusyBeaver/Research/ZFC/functional_subset_closure.md}{complete comparison}.
\bibitem{decoded} ProveIt source57, \emph{Two further source routes and an
explicit boundary}, section2,
\href{https://github.com/VladimirReshetnikov/ProveIt/blob/e58338db7545909907e07a7603e8eb88625c527a/SetTheory/Cardinals/docs/reports/set-theory/zfc-closure-kernels/57-prefix-logic_audit-further_semantic_routes.md}{decoded-cache MP and its fused implementation plan}.
\end{thebibliography}
\end{document}
